% ---------------------------------------------------------------------
%  Chapitre 3 — Juger les stratégies malgré la dépendance
% ---------------------------------------------------------------------
\chapter{Juger les stratégies malgré la dépendance : une escalade de prudence}
\label{chap:juger}

Ce chapitre traite la question : parmi les dix stratégies d'analyse
technique, auxquelles s'ajoute le témoin, lesquelles produisent un avantage
(\edgevar) significativement
positif ? On applique d'abord le test de référence sous ses hypothèses
standard, puis on met à l'épreuve, une à une, ses conditions
d'application — en s'appuyant sur la structure de dépendance mesurée au
chapitre~\ref{chap:donnees}. On terminera en interrogeant le marché
lui-même, pour savoir si le verdict obtenu était prévisible.

% =====================================================================
\section{Chaque stratégie bat-elle le hasard ?}
\label{sec:b1-etape1}

\subsection{La question posée à l'outil}

La première question posée à l'outil est donc : l'avantage moyen affiché
par une stratégie est-il réel, ou n'est-il qu'un accident
d'échantillonnage ? On la formalise en confrontant la moyenne de
l'\edgevar{} à la valeur zéro, ce qui conduit au test de Student présenté
ci-dessous.

\subsection{Les méthodes qui y répondent}

% --- Le test de Student ---
\subsubsection*{Test de Student}

Lorsque l'on obtient un \edgevar{} moyen positif, on cherche à savoir si la
stratégie constitue réellement une prédiction, ou si cela est dû à l'aléa
de l'échantillon.

Il faut ici distinguer deux quantités : l'\edgevar{} moyen
théorique \(\mu_{\edgevar}\), c'est-à-dire la vraie valeur (inconnue)
gouvernant la stratégie, et sa moyenne
observée \(\bar e\), calculée sur notre échantillon de trades
(connue). On se pose la question suivante : à partir de la valeur observée
\(\bar e\), peut-on conclure que la vraie valeur \(\mu_{\edgevar}\) est
positive ?

Le test de Student~\cite{student1908} à un échantillon répond à cette question en confrontant
\(\mu_{\edgevar}\) à une valeur de référence fixée à zéro (zéro~$=$ aucun
avantage sur le hasard). Sa mise en \oe uvre suppose deux conditions : que l'\edgevar{}
vérifie un critère de normalité, et que les trades soient indépendants les uns des autres.
La première sera assouplie par le TCL (ci-après) ; 
la seconde, plus délicate, sera examinée en section suivante. On oppose ainsi deux
hypothèses :
\[
  \begin{cases}
    \Hzero : \mu_{\edgevar} = 0 & \text{(le vrai \edgevar{} est nul : pas mieux que le hasard),}\\[2pt]
    \Hun : \mu_{\edgevar} > 0 & \text{(le vrai \edgevar{} est positif : avantage réel).}
  \end{cases}
\]
Le test est unilatéral : seul un avantage strictement positif nous
intéresse.

Pour trancher, on définit \(t\) comme étant le nombre d'erreurs-types
nécessaires pour faire basculer l'\edgevar{} moyen de zéro vers celui obtenu
dans notre échantillon :
\[
  t = \frac{\bar e}{s/\sqrt{n}} \sim t_{n-1} \quad\text{sous } \Hzero .
\]

Sous \Hzero{} (\edgevar{} nul), \(t\) suit une loi de Student à \(n-1\) degrés
de liberté. On évalue alors la p-value, c'est-à-dire la probabilité
d'obtenir un \(t\) au moins aussi grand que celui observé si \Hzero{}
était vraie.

Si la p-value est très faible, le hasard n'est plus une
explication crédible de l'\edgevar{} moyen positif observé.

On tranche alors sur le fait que l'\edgevar{} moyen théorique est positif, et que la
stratégie permet donc de faire mieux que le hasard.


On a dit précédemment que le test de Student nécessitait que l'\edgevar{} suive
une loi normale. On cherche à vérifier cette condition, dans un premier
temps avec le test de Shapiro--Wilk, puis dans un second temps à l'aide du
théorème central limite (TCL). En effet, le test de Shapiro--Wilk a peu de
chances de valider cette hypothèse : la taille conséquente de l'échantillon
déstabilise ce test, alors qu'elle renforce au contraire le TCL.

% --- Le test de normalite ---
\subsubsection*{Test de normalité de Shapiro--Wilk}
Le test de Shapiro--Wilk~\cite{shapiro1965} est un test de normalité qui consiste à comparer
deux mesures de dispersion : d'une part les écarts à la moyenne, d'autre
part l'ordre des valeurs. Ces deux mesures ne coïncident que si
l'hypothèse de normalité est vraie. On oppose ainsi deux hypothèses :
\[
  \begin{cases}
    \Hzero \;(W \simeq 1) : \text{l'\edgevar{} suit une loi normale,}\\[2pt]
    \Hun \;(W \ll 1) : \text{l'\edgevar{} ne suit pas une loi normale.}
  \end{cases}
\]
On notera que, contrairement au test de Student, c'est ici \Hzero{} qui
correspond à la propriété recherchée (la normalité) : une p-value faible
conduira donc à rejeter la normalité.

On pose alors la statistique \(W = b^2 / \mathrm{SCE}\) comme rapport de
ces deux mesures. Lorsque \(W\) est proche de 1, l'échantillon est
compatible avec une loi normale.

Pour trancher sur la normalité, on procède de façon similaire au test de
Student : on calcule la p-value d'obtenir notre \(W\) observé si la loi
était normale ; si cette probabilité est trop faible, alors on rejette
l'hypothèse de normalité.

Appliqué à nos données, le test rejette la normalité pour les onze
stratégies, comme le montre le tableau~\ref{tab:shapiro}. La condition du
test de Student n'est donc, en apparence, pas remplie. On remarque que même
les stratégies dont le \(W\) est proche de 1, comme \texttt{dt\_top}
($W = 0{,}95$), sont rejetées : à grand \(n\), le moindre écart à la normale
suffit à produire une p-value négligeable, ce qui justifie de recourir au
théorème central limite plutôt qu'à ce test pour valider l'emploi de Student.

\begin{table}[H]
\centering
\caption{Statistique \(W\) et p-value de Shapiro--Wilk par stratégie
(\edgevar{} ; sous-échantillon de 5\,000 trades lorsque \(n > 5\,000\)).}
\label{tab:shapiro}
\small
\begin{tabular}{@{}lrrr@{}}
\toprule
Stratégie & $n$ & $W$ & $p$ (Shapiro) \\
\midrule
\texttt{oracle} (témoin) & 6\,798 & $0{,}590$ & $6{,}8\times10^{-76}$ \\
\midrule
\texttt{db\_bottom}    & 8\,513 & $0{,}911$ & $2{,}7\times10^{-47}$ \\
\texttt{dt\_top}       & 7\,381 & $0{,}950$ & $3{,}2\times10^{-38}$ \\
\texttt{hs\_classic}   & 1\,065 & $0{,}886$ & $2{,}4\times10^{-27}$ \\
\texttt{hs\_inverse}   & 1\,206 & $0{,}963$ & $6{,}4\times10^{-17}$ \\
\texttt{ma\_crossover} & 5\,417 & $0{,}498$ & $3{,}2\times10^{-80}$ \\
\texttt{rsi\_classic}  & 7\,117 & $0{,}871$ & $1{,}5\times10^{-53}$ \\
\texttt{rsi\_strict}   &   978  & $0{,}717$ & $1{,}1\times10^{-37}$ \\
\texttt{rsi\_trend}    & 2\,789 & $0{,}883$ & $1{,}2\times10^{-41}$ \\
\texttt{sr\_breakdown} & 1\,377 & $0{,}775$ & $8{,}0\times10^{-40}$ \\
\texttt{sr\_breakout}  & 2\,367 & $0{,}843$ & $2{,}1\times10^{-43}$ \\
\bottomrule
\end{tabular}
\end{table}

% --- Le theoreme central limite ---
\subsubsection*{Théorème central limite}
Le rejet de Shapiro--Wilk porte sur la non-normalité de l'\edgevar{}
individuel. Mais comme le test de Student ne dépend que de la
moyenne de l'\edgevar{}, on peut se rabattre sur le théorème central
limite~\cite{saporta2006}, qui garantit la quasi-normalité de la moyenne sur
les grands échantillons, quelle que soit la loi de départ :
\[
  \frac{\bar e - \mu_{\edgevar}}{\sigma/\sqrt{n}}
  \;\xrightarrow[n \to \infty]{}\; \mathcal{N}(0,1) .
\]
Cette justification suppose un effectif suffisant ;
un outil générique devrait écarter les stratégies trop peu fournies.
Ici, \(n\) est compris entre 978 et 8\,513 selon la stratégie, ce qui est
largement suffisant.
La condition réellement utile au test de Student est donc satisfaite,
malgré le rejet de Shapiro--Wilk.

On remarque de plus que cette quantité est exactement celle du test de
Student : sous \Hzero{} (\(\mu_{\edgevar} = 0\)), et en remplaçant l'écart-type
théorique \(\sigma\) (inconnu) par son estimation \(s\), on retrouve la
statistique \(t = \bar e / (s/\sqrt{n})\). Le théorème central limite
justifie ainsi directement la forme et la validité du test de Student.

% --- La correction de Bonferroni ---
\subsubsection*{Correction de Bonferroni}
Pour appliquer le test de Student, il faut fixer une valeur seuil pour la
p-value. Lorsque l'on effectue un test isolé, on prend habituellement
\(0{,}05\) : on estime que s'il n'y avait que \(5\,\%\) de chances
d'obtenir un tel résultat sous \Hzero{}, il est peu vraisemblable d'être
tombé dessus par hasard, et l'on accepte alors le risque de rejeter
\Hzero{}.

Seulement, ici, on effectue \(11\) tests. Plaçons-nous dans le cas le plus
défavorable où aucune stratégie n'a d'avantage réel (\Hzero{} vraie pour
les onze) : pour un test isolé, la probabilité de ne pas conclure à tort à
un avantage est alors \(0{,}95\). En supposant les tests indépendants, la
probabilité de ne commettre aucune erreur sur les \(11\) tests tombe à
\(0{,}95^{11} \approx 0{,}57\). Le risque de commettre au moins un
faux positif sur l'ensemble s'élève donc à
\(1 - 0{,}95^{11} \approx 0{,}43\), ce qui est inacceptable.

Pour éviter cela, on raisonne en sens inverse : on fixe le risque
global à \(0{,}05\) et l'on en déduit le seuil de chaque test.
On s'appuie sur l'\textbf{inégalité de Boole}, selon laquelle la
probabilité qu'au moins un événement se produise est inférieure ou égale
à la somme de leurs probabilités individuelles. Pour garantir un risque
global inférieur à \(0{,}05\), il suffit donc que la somme des seuils
individuels soit égale à \(0{,}05\). En les prenant tous égaux, on obtient
la \textbf{correction de Bonferroni} :
\[
  \alpha' = \frac{\alpha}{m} = \frac{0{,}05}{11} \approx 0{,}0045 .
\]

\subsection{Résultats}
\begin{table}[H]
\centering
\caption{Test de l'avantage moyen par stratégie (seuil Bonferroni $\approx 0{,}0045$).}
\label{tab:b1-etape1}
\small
\begin{tabular}{@{}lrrrc@{}}
\toprule
Stratégie & $n$ & \edgevar{} moyen & $p$ (Student) & Verdict \\
\midrule
\texttt{oracle} (témoin) & 6\,798 & $+0{,}0658$ & $\approx 0$    & \textbf{bat le hasard} \\
\midrule
\texttt{rsi\_classic} & 7\,117 & $+0{,}0123$ & $1{,}2\times10^{-8}$  & \textbf{bat le hasard} \\
\texttt{rsi\_strict}  &   978  & $+0{,}1313$ & $5{,}5\times10^{-10}$ & \textbf{bat le hasard} \\
\texttt{rsi\_trend}   & 2\,789 & $+0{,}0021$ & $0{,}25$             & non \\
\texttt{db\_bottom}   & 8\,513 & $-0{,}0025$ & $0{,}98$             & non \\
\texttt{ma\_crossover}& 5\,417 & $-0{,}0248$ & $\approx 1$          & non \\
\multicolumn{5}{c}{\footnotesize (autres stratégies : $p>0{,}45$, non significatives)} \\
\bottomrule
\end{tabular}
\end{table}

Outre l'oracle témoin (qui bat le hasard de très loin, comme attendu),
seules \texttt{rsi\_classic} et \texttt{rsi\_strict} survivent à la
correction de Bonferroni.

Il est judicieux de rappeler que ce test vient uniquement répondre à la question de savoir 
si le \edgevar{} moyen positif est dû seulement à l'échantillon, auquel cas on conserve \Hzero{}, ce qui
signifie que la stratégie n'est pas gagnante.
Toutefois dans le cas inverse, on pourra affirmer que la stratégie a capturé un signal,
mais pas qu'elle est gagnante pour deux raisons. 
D'une part, un \edgevar{} positif ne signifie pas des gains positifs, mais simplement que l'on perd moins que le hasard. 
D'autre part, dû au fait que n est grand, la loi normale de la moyenne a une largeur qui s'écrase : 
on fait alors ressortir comme significatives des stratégies dont la moyenne
d'\edgevar{} n'est que très
faiblement positive, que les frais suffiraient à rendre perdantes.

Surtout, ce verdict repose encore sur une condition que nous n'avons pas
vérifiée, celle de l'indépendance des trades. C'est l'objet de la section
suivante.




% =====================================================================
\section{Le verdict survit-il à la dépendance des trades ?}
\label{sec:b1-dependance}

\subsection{La question posée à l'outil}
\label{sec:question-dependance}
On a appliqué de façon naïve le test de Student en supposant que
les données étaient indépendantes. Or, comme nous l'avons mesuré
au chapitre~\ref{chap:donnees}, nos données sont dépendantes par le titre,
par le marché et d'un mois à l'autre.
Ces dépendances empruntent en réalité deux chemins seulement, que nous
appellerons les deux portes. La première est la porte de l'action, celle des
trades ouverts sur un même titre. La seconde est la porte de la période,
celle des trades ouverts à une même époque, qu'ils subissent ensemble le
mouvement du marché ou qu'ils se prolongent sur les mêmes mois. Ce sont les
deux seules voies par lesquelles deux trades peuvent partager une cause
commune.

Le problème est que le test de Student traite alors les trades d'une stratégie
comme autant d'informations distinctes, alors que ces informations se
répètent en partie. L'erreur-type s'en trouve sous-estimée. Comme cette
erreur-type sert à mesurer l'écart entre l'avantage observé et la valeur de
référence, un dénominateur trop petit gonfle la statistique de Student et
rend la p-value trop optimiste. Le risque de conclure à tort qu'une stratégie
bat le hasard augmente d'autant.

Nous cherchons donc à savoir si le verdict positif obtenu pour
\texttt{rsi\_classic} et \texttt{rsi\_strict} résiste lorsque l'on abandonne
l'hypothèse d'indépendance. 
Pour cela, nous corrigeons d'abord l'erreur-type
au moyen de l'effet de sondage, puis nous vérifions ce résultat par un
rééchantillonnage qui ne suppose aucune forme particulière des données. Nous
fermons ensuite chaque porte en regroupant les trades pour que chaque titre,
puis chaque période, ne compte qu'une fois. Nous vérifions enfin que la porte
de la période tient elle-même, en traitant la suite des avantages mensuels
comme une série temporelle.

\subsection{Les méthodes qui y répondent}

% --- L'effet de sondage ---
\subsubsection*{Corriger l'erreur-type par l'effet de sondage}

Comme on l'a vu au chapitre~\ref{chap:donnees}, dans nos données, les titres
forment des grappes. On a également mesuré la corrélation intra-titre
$\hat\rho$, qui dit à quel point les trades d'un même titre se ressemblent.

On va désormais introduire l'effet de sondage (design effect) qui est le rapport
entre la variance d'un plan de sondage et la variance d'un plan de sondage de
référence, typiquement le sondage aléatoire simple.
Ici ce sera le rapport entre la variance réelle (les trades groupés par titre)
et la variance de référence (sondage aléatoire simple).

\[
  \mathrm{DEFF}
    = \frac{\mathrm{Var}(\hat\theta)}{\mathrm{Var}_{\mathrm{SAS}}(\hat\theta)} .
\]

Ici il s'agit d'un sondage à plusieurs degrés ; le cours de théorie des
sondages établit qu'alors l'effet de sondage vaut

\[
  \mathrm{DEFF} = 1 + (\bar m - 1)\,\hat\rho ,
\]

où $\bar m$ désigne le nombre moyen de trades par titre, et $\hat\rho$ la
corrélation intra-titre mesurée au chapitre~\ref{chap:donnees}.
La variance réelle de la moyenne vaut $\mathrm{DEFF}$ fois celle que l'on
obtiendrait sous indépendance. On doit alors considérer que l'on ne possède
seulement que $n_{\mathrm{eff}}$ informations.

\[
  n_{\mathrm{eff}} = \frac{n}{\mathrm{DEFF}} .
\]

Les largeurs d'IC sont alors multipliées par $\sqrt{\mathrm{DEFF}}$ par rapport au
plan de sondage de référence

\[
  \widehat{\mathrm{se}}_{\mathrm{corr}}
        = \frac{s}{\sqrt{n}}\sqrt{\mathrm{DEFF}} .
\]

On refait le test de Student de la section~\ref{sec:b1-etape1} avec cette
erreur-type corrigée, en
conservant le même seuil de Bonferroni.

% --- Le bootstrap par grappes ---
\subsubsection*{Vérifier sans hypothèse de forme}

La méthode ne va pas tester la dépendance intra-titre,
mais plutôt évaluer son impact sur le \edgevar{} moyen.
Au lieu d'utiliser le test de Student, on va directement, à partir des
grappes, essayer de créer une courbe de Gauss qui sera représentante de la
répartition du \edgevar{}, avec pour objectif de vérifier l'impact de la dépendance
sur le \edgevar{}. Car si la dépendance a un impact fort, alors notre cloche sera
large et la probabilité d'être en dessous de zéro va augmenter.

Pour produire cette courbe, on va, comme à la méthode précédente, définir
chaque titre comme étant une grappe à l'intérieur de laquelle les données
sont dépendantes. On effectue alors un grand nombre de tirages, ici 4\,000.
À chaque tirage, on tire avec remise autant de titres que la stratégie en a
mobilisés (soit environ cinq cents titres). Chaque titre peut être tiré
plusieurs fois, auquel cas on prend autant de fois l'ensemble des trades
qu'il contient. On calcule ensuite la moyenne des edges ainsi obtenue par
ce tirage.

On obtient ainsi, après les 4\,000 tirages, une cloche dont la moyenne ne
diffère pas de la moyenne initiale : le tirage étant uniforme sur les titres,
chaque trade a la même chance d'être repioché, si bien que la moyenne des
tirages se recentre sur celle de départ. En revanche, la largeur de la cloche
nous permet d'avoir une vision de l'impact de la dépendance intra-titre sur
notre échantillon.

Il est utile de préciser que l'on calcule la moyenne des edges de l'ensemble
des trades obtenus, afin de garder le poids qu'a chaque titre, contrairement
à la méthode de la porte que l'on verra plus loin. C'est justement en
gardant le poids de chaque titre que l'on conserve l'information de la
dépendance intra-titre et que l'on met en valeur son influence sur le \edgevar{}
moyen.

En effet, lorsque la dépendance est très impactante, on obtient une cloche
large qui va déborder sous zéro. On calcule alors le nombre de tirages qui
passent sous zéro, que l'on divise par le nombre total de tirages. Cette
proportion est la p-value du test : si elle dépasse le seuil défini par la
correction de Bonferroni, on rejette la stratégie.

Cette méthode suppose néanmoins que les titres soient indépendants entre
eux, ce qui n'est pas le cas ici à cause du facteur marché. Elle reste
toutefois un premier filtre légitime : en ignorant cette dépendance, le
bootstrap sous-estime la largeur de la cloche et se montre donc trop
permissif, si bien qu'une stratégie qui échoue déjà à ce stade est éliminée
sans appel, tandis que les autres restent à départager par la porte suivante.
Cette technique de rééchantillonnage est le bootstrap, proposé par
Efron~\cite{efron1979}.

Attention, cette méthode met en évidence le rôle de la dépendance intra-titre
sur le \edgevar{}. Cependant, le fait de conserver le poids de chaque titre nous laisse
vulnérables. En effet, si la grande majorité des trades sont obtenus seulement sur
un petit nombre de titres, alors nos tirages seront généralement similaires,
et donc notre courbe sera étroite.
C'est pourquoi, avec cette méthode, une cloche étroite ne suffit pas à
conclure que l'avantage est réel.
Pour lever ce doute, il faut retirer aux gros titres le poids que le nombre
de leurs trades leur donne. C'est l'objet de la méthode suivante.

% --- Les deux portes ---
\subsubsection*{Fermer chaque porte en donnant un vote à chaque unité}

Nous avons mesuré qu'il y a trois canaux de dépendance qui passent par deux portes.
La première porte, celle de l'action ; la seconde porte, celle de la période, qui
englobe la dépendance du marché ainsi que le chevauchement.

On utilise la méthode de l'agrégation, qui consiste à calculer l'avantage moyen
de chaque grappe puis à appliquer le test de Student sur ces données, avec une
seule voix par grappe quel que soit le nombre de trades qu'elle contient.
Pour la porte de l'action, on agrégera par titre. Pour la porte de la période,
on agrégera par mois : chaque trade est rattaché à son mois d'entrée, celui où
le signal s'est déclenché et où la décision a été prise. On évite ainsi que la
fluctuation du marché sur une même période, ou les trades qui se chevauchent,
ne soient comptés comme plusieurs informations distinctes.
Dans les deux cas, la dépendance à l'intérieur d'une grappe ne peut plus rien
fausser, puisqu'un titre, ou un mois, ne pèse qu'une seule voix. En revanche,
rattacher un trade à son seul mois d'entrée alors qu'il peut rester ouvert
plusieurs mois lie les mois voisins entre eux : cette dépendance inter-mensuelle
sera le sujet de la méthode suivante.
N'oublions pas que le théorème central limite s'applique alors à ces moyennes
de grappes, dont le nombre reste élevé : selon les stratégies, il reste 431 à
501~titres pour la porte de l'action et 138 à 197~mois pour celle de la période,
ce qui est bien suffisant. Le protocole ayant vocation à juger les stratégies à
venir, il signale toutefois le cas où une stratégie trop peu diffusée
laisserait un nombre de grappes insuffisant pour invoquer le théorème : le
verdict serait alors ininterprétable et non pas défavorable. Pour être validée,
une stratégie doit franchir les deux portes, chacune au seuil de Bonferroni.

Cette technique fait perdre beaucoup d'information. Mais on privilégiera
le risque de faux négatif à celui de faux positif.

Ayant déjà évalué la dépendance intra-titre par les méthodes DEFF et bootstrap,
l'apport de cette méthode sera de mettre en lumière les stratégies qui, comme
nous l'avons expliqué à la fin de la méthode d'Efron, reposent uniquement sur
quelques titres, car l'agrégation supprime le poids des titres.

On rejettera les stratégies qui ne passeront pas la fermeture des portes.
Toutefois, une stratégie rejetée parce qu'elle dépend de certains titres en
particulier pourra être réétudiée afin de peaufiner la sélection des trades.
Il est possible qu'une stratégie ne soit efficace que sur un secteur en
particulier. Il faudra évidemment faire attention au sur-apprentissage.

Une objection mérite d'être levée. La porte de l'action suppose les titres
indépendants entre eux, ce qui est contestable puisque le marché explique une
grande part de la variance : deux titres tradés au même moment subissent le
même choc. Mais cette dépendance-là est de nature temporelle, deux titres ne
sont liés que s'ils sont tradés à la même période, et c'est précisément ce que
la porte de la période traite. Les deux portes se complètent donc, chacune
prenant en charge une source, ce qui justifie de les poser toutes les deux.

Elles ne sont toutefois pas fermées simultanément : nous conduisons deux tests
séparés plutôt qu'une agrégation croisée par titre et par mois, qui produirait
des grappes bien plus petites et rouvrirait à moitié chacune des deux portes.
La règle du ET est d'ailleurs plus exigeante, puisqu'une stratégie doit
survivre à chaque source de dépendance prise isolément et traitée dans le pire
des cas. Seule échapperait une dépendance propre à l'intersection, un effet
qui n'existerait que pour un titre donné pendant un mois donné, et l'on voit
mal comment un tel effet porterait à lui seul un avantage. Le traitement
simultané des deux dimensions existe, c'est le double clustering de
Petersen~\cite{petersen2009}, mais il sort du cadre des méthodes simples
retenues ici ; nous y revenons en section~\ref{sec:limites}.

% --- La serie mensuelle ---
\subsubsection*{Vérifier que la porte de la période tient elle-même}

Comme nous l'avons mesuré au deuxième canal du chapitre~\ref{chap:donnees}, le
marché explique 73~\% de la variance : lorsqu'un événement survient, tout le
marché réagit de façon similaire. L'agrégation mensuelle regroupe ainsi les
trades déclenchés par un même contexte de marché, et rattacher chaque trade à
son seul mois d'entrée garantit qu'aucun ne pèse dans deux grappes à la fois.

Ce rattachement règle le double comptage, mais non le recouvrement des périodes
de détention. Nous avons vu au chapitre~\ref{chap:donnees} que les trades se
prolongent souvent au-delà d'un mois : un trade ouvert en janvier et clos en
avril est rangé dans janvier, alors que son résultat dépend de ce qui s'est
passé jusqu'en avril. La grappe de janvier porte donc de l'information sur les
mois suivants, que la grappe de février porte également. Aucun trade n'est
compté deux fois, et pourtant deux mois voisins partagent les mêmes conditions
de marché ; avec 73~\% de la variance expliquée par ces conditions communes, ce
recouvrement suffit à les faire se ressembler.

La suite des avantages mensuels est donc autocorrélée, et le test de Student
avec agrégation mensuelle pourrait être trop optimiste.
On voudrait alors appliquer comme précédemment la méthode du DEFF pour corriger
l'erreur-type. Mais nos données sont chronologiques : la dépendance y est
ordonnée, un mois étant plus lié au précédent qu'à celui d'il y a cinq mois,
là où les trades d'un même titre étaient interchangeables. Un coefficient
unique ne peut donc pas la décrire, et il nous faut une autre méthode.

Dans cette méthode nous voulons calculer, pour un mois~$t$, l'impact moyen
qu'ont les mois qui le précèdent, en cherchant jusqu'à six mois en arrière.
Pour cela nous allons traiter cette suite comme une série temporelle. On
supposera qu'il n'y a pas d'effet de saison, et surtout que le lien entre deux
mois ne dépend que de leur écart et non de leur date, le mois de janvier
réagissant comme les autres mois de l'année.

On mesure alors l'impact propre de chaque décalage par l'autocorrélation
partielle (PACF), qui retranche l'influence des mois intermédiaires ; un
mois peut en effet sembler lié à celui d'il y a deux mois par le seul fait
qu'ils sont tous deux liés au mois qui les sépare. Cette autocorrélation
partielle s'annule au-delà du dernier mois qui compte vraiment, ce qui nous
donne le nombre de mois à retenir, puis l'impact de chacun d'eux. Ensuite, de
la même façon que pour la méthode du DEFF, on corrigera l'erreur-type du test
de Student avec agrégation mensuelle.

La méthode se déroule alors en trois temps, chacun porté par un outil :
y a-t-il réellement une dépendance (test de Box et Pierce) ? peut-on la
modéliser (stationnarité par Dickey--Fuller et KPSS, ordre par
l'autocorrélation partielle, poids par Yule-Walker) ? de combien corriger
(variance de long terme) ? Les tests de conditions qui jalonnent ce parcours
sont menés au seuil usuel de 5~\% ; le seuil de Bonferroni reste réservé au
verdict final.

% --- Le test de Box et Pierce ---
\subsubsection*{Test d'indépendance de Box et Pierce}

Avant de corriger, encore faut-il qu'il y ait quelque chose à corriger. Le
raisonnement du recouvrement rend l'autocorrélation plausible, il ne la rend
pas certaine, et elle peut différer d'une stratégie à l'autre. Le protocole
commence donc par poser directement la question à chaque série : les mois
sont-ils indépendants ? Le test de Box et Pierce~\cite{boxpierce1970} y répond sur la série brute,
sa statistique cumule les six premières autocorrélations, mises au carré pour
que les positives et les négatives ne se compensent pas, et mises à l'échelle
de la longueur de la série :
\[
  S_{\mathrm{BP}} \;=\; T \sum_{h=1}^{6} \hat\rho(h)^2 ,
\]
où \(T\) est le nombre de mois de la série et \(\hat\rho(h)\)
l'autocorrélation d'ordre \(h\), c'est-à-dire la corrélation entre la série
et elle-même décalée de \(h\) mois.

Ce nombre de six décalages n'est pas choisi au hasard. La dépendance que nous
cherchons vient du recouvrement des périodes de détention, un trade ouvert en
janvier et clos en avril liant janvier aux mois suivants. Elle ne peut donc
pas s'étendre au-delà de la durée des trades, qui dépassent rarement six mois.
Regarder plus loin reviendrait à chercher un effet là où aucune cause ne peut
le produire, tout en ajoutant du bruit à la statistique. Ce même plafond nous
servira plus loin pour choisir l'ordre du modèle, et nous vérifierons en fin
de méthode qu'aucune mémoire ne subsiste au-delà.

On oppose alors :
\[
  \begin{cases}
    \Hzero \;(S_{\mathrm{BP}} \text{ petit}) : \text{les mois sont indépendants,}\\[2pt]
    \Hun \;(S_{\mathrm{BP}} \text{ grand}) : \text{au moins un décalage porte une autocorrélation.}
  \end{cases}
\]
C'est \Hzero{} qui porte ici la propriété recherchée.

Sous \Hzero{}, chaque $\hat\rho(h)$ n'est que du bruit d'échantillonnage
d'ordre $1/\sqrt{T}$ : la somme $S_{\mathrm{BP}}$ suit alors une loi du
khi-deux à six degrés de liberté. On calcule donc la p-value
de notre $S_{\mathrm{BP}}$ observé ; si elle est trop faible, l'indépendance
n'est plus une explication crédible.

Si l'indépendance n'est pas rejetée, la condition du test par
agrégation mensuelle est remplie telle quelle : son verdict tient, et aucune
correction n'est requise — on ne corrige que ce qui est mesuré, comme pour le
DEFF. Si elle est rejetée, l'autocorrélation est réelle et la suite de la
méthode s'applique.

Pour effectuer cette correction on va supposer que les mois sont tous
similairement exposés aux mois qui les précèdent ; pourtant cette affirmation
n'est vraie que si la série est stationnaire, c'est pourquoi on va effectuer
deux tests pour vérifier la stationnarité de la série de chaque stratégie
concernée.

Si une série ne l'était pas, on sortirait du domaine de validité de la
méthode. Le protocole la signale et s'abstient de conclure plutôt que de
produire un verdict fondé sur un impact moyen qui n'aurait pas de sens. Aucune
de nos stratégies n'est dans ce cas, mais la situation peut se présenter pour
une stratégie à venir.

% --- Stationnarité, versant 1 -----------------------------------------
\subsubsection*{Test de stationnarité de Dickey et Fuller augmenté}

Ce test, comme le suivant, se calcule directement sur la série.
Il consiste à calculer, par une simple régression du
mois courant sur le mois précédent, un coefficient que nous noterons \(\phi\) :
\(x_t = \phi\, x_{t-1} + \eta_t\).
En remplaçant \(x_{t-1}\) par sa propre expression, puis en recommençant, on
voit ce que devient un choc au fil du temps :
\[
  x_t = \eta_t + \phi\,\eta_{t-1} + \phi^2\,\eta_{t-2} + \phi^3\,\eta_{t-3} + \cdots
\]
Le mois courant est donc la somme de tous les chocs passés, chacun pesé par
\(\phi\) élevé à son ancienneté. Si \(\phi\) est inférieur à 1, ces puissances
tendent vers zéro : l'impact d'un mois sur ceux qui suivent décroît, la série
oublie ses chocs et revient vers sa moyenne, elle est stationnaire. Si \(\phi\)
vaut 1, toutes les puissances valent 1 : chaque choc s'ajoute définitivement au
niveau, la série dérive sans repère et n'est donc pas stationnaire.
On oppose ainsi :

\[
  \begin{cases}
    \Hzero \;(\phi = 1) : \text{les chocs s'accumulent, la série n'est pas stationnaire,}\\[2pt]
    \Hun \;(\phi < 1) : \text{les chocs s'estompent, la série est stationnaire.}
  \end{cases}
\]

Le test est unilatéral : un \(\phi\) supérieur à 1 décrirait une série
explosive, cas sans objet ici. C'est donc le rejet qui conclut à la
stationnarité.

Comme pour le test de Student, la statistique est le nombre d'erreurs-types
qui séparent le \(\hat\phi\) estimé de la valeur 1 :

\[
  t_{\mathrm{DF}} = \frac{\hat\phi - 1}{\widehat{\mathrm{se}}(\hat\phi)},
\]


où \(\hat\phi\) est le coefficient estimé par la régression et
\(\widehat{\mathrm{se}}(\hat\phi)\) son erreur-type, c'est-à-dire
l'incertitude qui entoure cette estimation. Plus \(t_{\mathrm{DF}}\) est
négatif, plus \(\hat\phi\) s'écarte de 1 vers le bas et plus \Hzero{} devient
implausible.

Cette statistique, que les auteurs eux-mêmes nomment rapport de Student, n'en
suit pourtant pas la loi~\cite{dickeyfuller1979}. La raison tient à \Hzero{} :
sous cette hypothèse la série n'est justement pas stationnaire, sa variance
croît avec le temps, et les conditions qui donnent la loi de Student ne sont
plus réunies. Sa loi est alors celle que Dickey et Fuller ont établie et
tabulée, et qui porte leur nom. La lecture reste ensuite la même :
p-value confrontée au seuil, rejet si elle est trop faible.

Ce test suppose que les résidus de la régression ne soient
plus autocorrélés ; on emploie pour cela sa version dite augmentée, qui
ajoute des mois passés supplémentaires dans la régression jusqu'à les
nettoyer. Reste à savoir combien : trop peu et la mémoire n'est pas absorbée,
trop et chaque coefficient supplémentaire, estimé sur les mêmes données,
ajoute du bruit. C'est le critère d'information d'Akaike~\cite{akaike1974}
qui arbitre :
\[
  \mathrm{AIC}_k = 2\ln \hat\sigma_k + \frac{k}{T},
\]
où \(\hat\sigma_k\) mesure l'erreur résiduelle du modèle à \(k\) retards.
Le premier terme diminue quand le modèle colle mieux aux données, le second
augmente avec le nombre de coefficients : on retient le \(k\) qui minimise
la somme, de sorte qu'un retard n'est conservé que s'il apporte plus qu'il
ne coûte. Sur nos séries, ce nombre va de zéro à neuf selon les stratégies :
la condition est donc prise en charge par le test lui-même.

% --- Stationnarité, versant 2 -----------------------------------------
\subsubsection*{Test de stationnarité KPSS}

À l'inverse, le KPSS~\cite{kpss1992} cumule les écarts à la moyenne au fil du temps : si la
série oscille autour d'un niveau fixe, ces cumuls se compensent et restent
bornés, tandis qu'ils s'emballent dès que la série dérive. Ses hypothèses
sont inversées :
\[
  \begin{cases}
    \Hzero \;(\text{cumuls bornés}) : \text{la série est stationnaire,}\\[2pt]
    \Hun \;(\text{cumuls qui s'emballent}) : \text{la série n'est pas stationnaire.}
  \end{cases}
\]
Comme pour le test de Shapiro--Wilk, c'est ici \Hzero{} qui porte la propriété
recherchée : c'est l'absence de rejet qui conclut à la stationnarité.

Sa statistique somme les carrés de ces cumuls, rapportés à la longueur de la
série et à sa variance de long terme, de sorte qu'elle reste petite quand les
cumuls sont bornés et explose quand ils dérivent. Sa loi sous \Hzero{} est,
là aussi, tabulée par les auteurs : on rejette la stationnarité quand la
statistique dépasse la valeur critique, c'est-à-dire quand la p-value passe
sous le seuil.

Ce test demande de préciser autour de quoi la stationnarité est jugée ; nous
la testons autour d'une constante et non d'une tendance, car rien ne laisse
attendre qu'un avantage dérive régulièrement au fil des années.

% --- La règle de décision sur la stationnarité ------------------------
\subsubsection*{Concordance des deux tests}

Les deux hypothèses nulles étant opposées, chacun des tests protège l'autre de
sa propre faiblesse : ne pas rejeter n'a jamais valeur de preuve, et un test
seul ne distingue pas une série réellement stationnaire d'une série sur
laquelle il manque de puissance. Nous ne conclurons donc que lorsque les deux
concordent ; s'ils divergent, le diagnostic sera tenu pour non concluant
plutôt que pour défavorable.

Ces deux tests ne demandent aucun modèle, et c'est bien ainsi qu'il faut
l'entendre : le \(\phi\) de Dickey et Fuller ne sert qu'à trancher sur la
racine unitaire et n'est pas conservé. Les impacts que nous recherchons sont
d'une autre nature ; la stationnarité maintenant acquise nous donne le droit
de les estimer.

% --- L'ordre du modèle : combien de mois comptent ---------------------
\subsubsection*{Autocorrélation partielle et choix de l'ordre}

On mesure l'impact propre de chaque décalage \(h\) par l'autocorrélation
partielle \(\hat\tau(h)\). Reste à décider quels décalages comptent vraiment.
On procède décalage par décalage, en opposant :
\[
  \begin{cases}
    \Hzero \;\bigl(\tau(h) = 0\bigr) : \text{le mois } t-h \text{ n'a pas d'impact propre,}\\[2pt]
    \Hun \;\bigl(\tau(h) \neq 0\bigr) : \text{il en a un.}
  \end{cases}
\]
Sous \Hzero{}, \(\hat\tau(h)\) n'est que du bruit d'échantillonnage,
approximativement normal et d'écart-type \(1/\sqrt{T}\) : la statistique
\(\sqrt{T}\,\hat\tau(h)\) suit donc une loi normale centrée réduite. On en
tire, comme pour les tests précédents, la p-value de notre \(\hat\tau(h)\)
observé ; si elle passe sous le seuil de 5~\%, l'absence d'impact n'est plus
une explication crédible et le décalage est retenu. Cela revient à comparer
\(\bigl|\hat\tau(h)\bigr|\) à \(q_{0{,}975}/\sqrt{T}\), avec
\(q_{0{,}975} \approx 1{,}96\) — c'est la bande de significativité que l'on
trace sur les corrélogrammes.
L'ordre retenu \(\hat p\) est le dernier décalage significatif ; si
des décalages intermédiaires ne le sont pas, ils sont conservés
malgré tout, car on préfère surestimer la mémoire que la tronquer.

La recherche est plafonnée à six mois, pour la raison déjà donnée au test de
Box et Pierce. Cet argument reste toutefois une plausibilité, non une preuve.
Nous n'en resterons pas là : la validation du modèle, présentée plus bas,
vérifiera qu'aucune mémoire ne subsiste au-delà de ce plafond.

% --- L'estimation des poids -------------------------------------------
\subsubsection*{Équations de Yule-Walker}

L'ordre connu, on écrit le modèle autorégressif~\cite{boxjenkins1970} sur la
série centrée : le
mois courant est une somme pondérée des \(\hat p\) mois précédents, plus un
bruit,
\[
  x_t = \Phi_1\, x_{t-1} + \cdots + \Phi_{\hat p}\, x_{t-\hat p} + \eta_t ,
\]
où \(x_t\) est l'avantage moyen du mois \(t\) (écarté de la moyenne de la
série), \(\Phi_i\) le poids du mois \(t-i\) — l'impact propre que nous
cherchions —, et \(\eta_t\) le bruit, c'est-à-dire la part du mois courant
qu'aucun mois passé n'explique.
Ces poids ne se lisent pas directement dans les données. Les autocorrélations,
en revanche, se mesurent sans difficulté : il suffit de comparer la série à
elle-même décalée. Les équations de Yule-Walker font le lien entre les deux.
On les obtient en multipliant le modèle par \(x_{t-h}\), le mois situé \(h\)
crans en arrière, puis en prenant l'espérance : le bruit \(\eta_t\) disparaît,
étant indépendant de tout ce qui précède, et il reste, pour chaque
\(h = 1, \ldots, \hat p\) :
\[
  \rho(h) = \Phi_1\,\rho(h-1) + \Phi_2\,\rho(h-2) + \cdots + \Phi_{\hat p}\,\rho(h-\hat p),
\]
avec \(\rho(0) = 1\) et \(\rho(-k) = \rho(k)\), une série étant liée à
elle-même de la même façon dans les deux sens. Chaque autocorrélation
s'exprime ainsi en fonction des poids recherchés. Comme il y a autant
d'équations que de poids inconnus, on y reporte les autocorrélations mesurées,
on résout, et l'on remonte ainsi aux poids. Le cas d'un seul mois de mémoire
est éclairant : le système se réduit à \(\rho(1) = \Phi_1\), le poids est
alors exactement l'autocorrélation mesurée.

Une équation supplémentaire, celle du décalage nul, donne l'amplitude du
bruit :
\[
  \sigma^2_\eta = \gamma(0)\Bigl(1 - \sum_{i=1}^{\hat p} \Phi_i\,\rho(i)\Bigr),
\]
soit la variance de la série moins la part qu'expliquent les mois passés,
autrement dit ce qui reste imprévisible.

Ce système n'a de sens que si la série est stationnaire. Parler de
« l'autocorrélation à deux mois d'écart », c'est en effet supposer qu'il en
existe une seule valeur, la même en 2010 qu'en 2020. Sans cette garantie, il
en faudrait une par époque, et le système n'aurait plus de solution unique.

C'est la somme de ces poids qui nous intéresse, car elle mesure la persistance
totale de la série : c'est elle qui dira de combien l'erreur-type du test
mensuel doit être corrigée.

% --- La correction et le verdict --------------------------------------
\subsubsection*{Variance de long terme et test de Student corrigé}

Le modèle fournit alors une variance de la moyenne qui tient compte de la
persistance, appelée variance de long terme :
\[
  \gamma_{\mathrm{lr}}
     = \frac{\sigma^2_\eta}{\bigl(1-\sum_{i=1}^{\hat p}\Phi_i\bigr)^{2}},
  \qquad
  \widehat{\mathrm{Var}}(\bar x) \approx \frac{\gamma_{\mathrm{lr}}}{T},
  \qquad
  \mathrm{DEFF}_t = \frac{\gamma_{\mathrm{lr}}}{\gamma(0)},
\]
où \(\sigma^2_\eta\) et les \(\Phi_i\) sont le bruit et les poids estimés
ci-dessus, \(T\) le nombre de mois et \(\gamma(0)\) la variance de la série.
Cette variance de long
terme additionne toutes les autocovariances, là où la formule usuelle en
\(s^2/T\) suppose au contraire des mois sans lien : si les poids sont nuls,
le dénominateur vaut 1 et l'on retrouve la formule usuelle ; plus la
persistance est forte, plus la correction est importante. Le rapport
\(\mathrm{DEFF}_t\) est l'exact analogue temporel de l'effet de sondage de la
première méthode : un facteur qui dit combien de fois la variance de la
moyenne est sous-estimée quand on ignore la dépendance — sur l'axe du temps
cette fois, et obtenu par un modèle car la dépendance y est ordonnée.

Le test de Student de la porte de la période est alors repris tel quel, avec
deux changements : l'erreur-type devient
\(\sqrt{\gamma_{\mathrm{lr}}/T}\), et le nombre de degrés de liberté est
ramené à la taille effective \(T_{\mathrm{eff}} = T/\mathrm{DEFF}_t\), comme
pour le DEFF. La statistique diminue mécaniquement : une stratégie peut
franchir la porte de la période et échouer ici. Le verdict se lit au même
seuil de Bonferroni que les étapes précédentes.

% --- La validation finale ---------------------------------------------
\subsubsection*{Validation du modèle : Box et Pierce sur les résidus}

Reste à vérifier que le modèle a bien capté toute l'autocorrélation, sans
quoi la correction serait fondée sur une mémoire mal décrite. On dispose pour
cela des résidus \(\hat\eta_t\), c'est-à-dire ce que le modèle
n'explique pas : pour chaque mois, la valeur observée moins celle que les
mois précédents laissaient attendre. Le raisonnement est alors le même qu'au
début de la méthode, mais appliqué à ce qui reste : si le modèle a bien
absorbé toute la mémoire, les résidus n'en contiennent plus, et ils sont
indépendants d'un mois à l'autre. On leur pose donc exactement la question
qu'on avait posée à la série brute, avec le même test :
\[
  \begin{cases}
    \Hzero \;\bigl(S_{\mathrm{BP}} \text{ petit}\bigr) : \text{les résidus sont indépendants, le modèle a tout absorbé,}\\[2pt]
    \Hun \;\bigl(S_{\mathrm{BP}} \text{ grand}\bigr) : \text{il subsiste de la mémoire, le modèle est inadéquat.}
  \end{cases}
\]
La démarche est donc circulaire au bon sens du terme : le premier test
demandait y a-t-il de la mémoire ?, celui-ci demande en
reste-t-il ? Deux différences seulement. L'inspection porte sur quinze
décalages et non plus six, comme le recommande la pratique
usuelle~\cite{boxpierce1970} qui situe ce nombre entre quinze et vingt ; le
plafond de six mois n'a donc rien pu laisser échapper. Et comme \(\hat p\)
coefficients ont été estimés sur ces mêmes données, ce qui rend les résidus
artificiellement un peu plus propres, on retranche le nombre de paramètres du
modèle : la statistique suit sous \Hzero{} une loi du khi-deux à
\(15 - \hat p\) degrés de liberté.

Ne pas rejeter valide donc la correction. Un rejet signalerait au contraire
un ordre trop faible ou une forme inadaptée ; le protocole le signale et
s'abstient alors de conclure, comme pour la stationnarité.

Au terme de cette méthode, la condition d'indépendance entre mois, que la
porte de la période supposait sans la vérifier, a donc été soit établie, soit
prise en charge par une correction elle-même contrôlée. Rappelons que les
conditions ne s'appliquent qu'à la branche qui les appelle : une stratégie
dont l'indépendance mensuelle n'est pas rejetée conserve le verdict de la
porte de la période sans qu'on ait à exiger d'elle ni stationnarité ni modèle,
puisqu'aucune correction ne lui est appliquée.

\subsection{Résultats}

On a vu précédemment que seules trois stratégies passaient le premier filtre
du test de Student naïf, qui supposait l'indépendance. Analysons les résultats
du protocole sur ces stratégies.

% --- Resultats : DEFF et bootstrap ---
\begin{table}[H]
\centering
\caption{Diagnostic de dépendance intra-titre et test corrigé par l'effet de
sondage, avec vérification par le bootstrap.}
\label{tab:b1-dep-deff}
\small
\begin{tabular}{@{}lrrrccc@{}}
\toprule
Stratégie & $\hat\rho$ & DEFF & $n_{\text{eff}}$ & $p$ naïve & $p$ corrigée & $p$ bootstrap \\
\midrule
\texttt{oracle} (témoin) & $0{,}053$ & $1{,}67$ & $4\,076$ & $\approx 0$ & $\approx 0$ & $<2{,}5\times10^{-4}$ \\
\midrule
\texttt{rsi\_classic} & $0{,}000$ & $1{,}00$ & $7\,117$ & $1{,}2\times10^{-8}$ & $1{,}2\times10^{-8}$ & $<2{,}5\times10^{-4}$ \\
\texttt{rsi\_strict}  & $0{,}134$ & $1{,}15$ & $848$    & $5{,}5\times10^{-10}$ & $7\times10^{-9}$     & $<2{,}5\times10^{-4}$ \\
\multicolumn{7}{c}{\footnotesize (autres stratégies : verdict « non » inchangé, $\hat\rho \le 0{,}09$)} \\
\bottomrule
\end{tabular}
\end{table}

Le $\hat\rho$ intra-titre mesuré au chapitre~\ref{chap:donnees} était quasi nul
pour l'ensemble des stratégies, à l'exception de \texttt{rsi\_strict}
($0{,}134$). La correction DEFF, qui gonfle la variance en proportion de
cette dépendance, ne change donc presque rien : \texttt{rsi\_classic} et
\texttt{rsi\_strict} survivent toutes deux à cette première correction, et le
bootstrap par grappes, qui ne fait aucune hypothèse sur la forme de la
dépendance, confirme le verdict. Le témoin, lui aussi, franchit cette étape
sans difficulté.

On continue alors le protocole sur ces trois stratégies.

% --- Resultats : les deux portes ---
\begin{table}[H]
\centering
\caption{Le test des deux portes : une unité réelle, un vote.}
\label{tab:b1-dep-portes}
\small
\begin{tabular}{@{}lrrrrc@{}}
\toprule
 & \multicolumn{2}{c}{Porte action ($k$ titres)} & \multicolumn{2}{c}{Porte période ($k$ mois)} & \\
\cmidrule(lr){2-3}\cmidrule(lr){4-5}
Stratégie & $k$ & $p_A$ & $k$ & $p_B$ & Verdict \\
\midrule
\texttt{oracle} (témoin) & 501 & $\approx 0$          & 193 & $\approx 0$ & \textbf{OUI} \\
\midrule
\texttt{rsi\_classic}  & 501 & $0{,}059$            & 195 & $0{,}039$ & non \\
\texttt{rsi\_strict}   & 456 & $3{,}9\times10^{-4}$ & 138 & $0{,}26$  & non \\
\texttt{rsi\_trend}    & 495 & $0{,}22$             & 175 & $0{,}29$  & non \\
\texttt{sr\_breakdown} & 457 & $0{,}37$             & 184 & $0{,}94$  & non \\
\multicolumn{6}{c}{\footnotesize (six autres stratégies : $p_A \ge 0{,}74$, toutes « non »)} \\
\bottomrule
\end{tabular}
\end{table}
Ce tableau change la lecture du chapitre. L'oracle passe les deux
portes de très loin ($p_A$ et $p_B$ pratiquement nuls) : son avantage,
construit pour être réel, se retrouve aussi bien titre par titre que mois par
mois. C'est le résultat attendu, mais il est important, car il montre que
l'outil ne rejette pas tout par excès de prudence et qu'il sait dire
« oui » quand l'avantage est effectivement là.

\texttt{rsi\_classic}, en revanche, tombe aux deux portes
($p_A = 0{,}059$ et $p_B = 0{,}039$, loin du seuil de Bonferroni
$0{,}0045$). Les 7\,117 trades de cette stratégie ne
constituaient pas 7\,117 témoignages indépendants, mais un nombre bien plus
restreint de titres et de mois, chacun compté plusieurs fois. Dès que
chaque titre, puis chaque mois, ne vote plus qu'une seule fois, l'avantage
s'efface.

\texttt{rsi\_strict} a un comportement différent : elle passe brillamment la
porte de l'action ($p_A = 3{,}9\times10^{-4}$ : son avantage est bien partagé
par un grand nombre de titres), mais échoue à la porte de la période
($p_B = 0{,}26$). Son avantage est réel titre par titre, mais
épisodique dans le temps, sans régularité mois après mois
démontrable. En l'état, il n'est donc pas exploitable.

Le protocole écarte ces deux stratégies, mais avec une note signalant qu'elles
présentent des résultats remarquables sur certains titres ou certains mois.
Il faudra se garder du sur-apprentissage, mais peut-être existe-t-il un
facteur qui expliquerait ces résultats. En l'état, on ne peut pas les laisser
passer.

En dehors du témoin, plus aucune stratégie ne franchit les deux portes. Cela
n'a rien d'étonnant : on opère sur les cinq cents plus grandes capitalisations
des États-Unis et, comme ces actions ont une forte tendance à croître, il faut
un résultat extrême pour dépasser le seuil de l'improbable.



% --- Resultats : la serie mensuelle ---
\begin{table}[H]
\centering
\caption{Série mensuelle des \edgevar{} : autocorrélation et test corrigé de
la variance de long terme.}
\label{tab:b1-dep-mensuel}
\small
\begin{tabular}{@{}lrrrrrc@{}}
\toprule
Stratégie & $T$ & $\hat\rho(1)$ & $\hat p$ (AR) & $\mathrm{DEFF}_t$ & $p$ corrigée & $p$ Box--Pierce \\
\midrule
\texttt{oracle} (témoin) & 193 & $0{,}18$ & 0 & $1{,}0$ & $9{,}4\times10^{-85}$ & $0{,}16$ \\
\midrule
\texttt{ma\_crossover} & 188 & $0{,}56$ & 3 & $6{,}4$ & $0{,}94$ & $0{,}14$ \\
\texttt{dt\_top}       & 197 & $0{,}38$ & 1 & $2{,}2$ & $0{,}97$ & $0{,}69$ \\
\texttt{db\_bottom}    & 197 & $0{,}36$ & 1 & $2{,}1$ & $0{,}82$ & $0{,}90$ \\
\texttt{hs\_inverse}   & 186 & $0{,}31$ & 1 & $1{,}9$ & $0{,}91$ & $0{,}93$ \\
\texttt{rsi\_classic}  & 195 & $0{,}23$ & 1 & $1{,}6$ & $0{,}08$ & $0{,}93$ \\
\texttt{rsi\_strict}   & 138 & $0{,}03$ & 0 & $1{,}0$ & $0{,}26$ & $0{,}28$ \\
\multicolumn{7}{c}{\footnotesize (quatre autres stratégies : $\mathrm{DEFF}_t \le 2{,}4$, aucune significative)} \\
\bottomrule
\end{tabular}
\end{table}



Cette dernière étape ne s'applique qu'aux stratégies encore en lice, celles
qui auraient franchi les deux portes. Pour chacune, on se demande d'abord si
les mois voisins sont réellement liés : le test de Box et Pierce sur la série
brute tranche entre une série sans mémoire, qui conserve alors le verdict de
la porte de la période sans correction, et une série autocorrélée, qui appelle
une correction.

Le tableau montre que la moitié des séries seulement sont autocorrélées et
suivent la branche corrigée, l'autocorrélation d'ordre un atteignant $0{,}56$
pour \texttt{ma\_crossover}. Les autres, dont l'oracle et \texttt{rsi\_strict},
ne rejettent pas l'indépendance des mois : leur verdict de la porte de la
période tient tel quel, sans correction.

Pour les séries autocorrélées, la correction n'est légitime que sous
stationnarité. Le croisement des deux tests l'assure ici pour toutes : Dickey
et Fuller rejette la racine unitaire ($p \le 0{,}005$), tandis que le KPSS ne
rejette pas la stationnarité. Après correction, aucun verdict ne bascule : la
variance mensuelle était bien sous-estimée, d'un facteur allant jusqu'à
$6{,}4$ pour \texttt{ma\_crossover}, mais aucune stratégie ne devient
significative pour autant. Comme attendu, l'oracle franchit lui aussi cette
dernière étape.



\section{Lecture du verdict : le marché est-il prévisible ?}
\label{sec:b1-efficience}

Aucune stratégie ne passe, et on peut se demander si c'est normal. Toutes
celles que nous avons testées cherchent la même chose, un signal dans le passé
du marché qui annoncerait la suite. Il est donc légitime de poser la question
directement au marché lui-même, sans passer par les stratégies. Nous
utilisons pour cela les rendements quotidiens introduits au
paragraphe~\ref{sec:donnees-secteurs}, en prenant la colonne du marché
agrégé, soit 4\,132 jours de cotation.

On applique le test de Box et Pierce à cette série, exactement comme on l'a
fait sur les séries mensuelles. Le résultat est net, l'indépendance est
rejetée très largement ($S_{\mathrm{BP}} = 93{,}7$ sur six retards, soit une
p-value indiscernable de zéro). Il y a donc bien de la mémoire dans le marché,
et on pourrait croire que cela contredit tout ce qui précède.

C'est ici qu'il faut regarder l'amplitude et non la seule p-value.
L'autocorrélation d'ordre un vaut $-0{,}08$, ce qui veut dire que la veille
explique $0{,}6\,\%$ de ce qui se passe le lendemain. Le signe négatif indique
un léger effet de rebond, une baisse ayant tendance à être suivie d'une
hausse. Si l'on construit la stratégie qui exploite ce rebond, en pariant
chaque jour contre le mouvement de la veille, elle rapporte $0{,}007\,\%$ par
jour avant frais. Or le moindre aller-retour en coûte davantage, si bien que
la stratégie devient perdante dès que l'on tient compte des frais les plus
faibles.

Le marché est donc prévisible au sens statistique du terme, mais il ne l'est
pas au sens où l'on pourrait en tirer de l'argent. C'est exactement ce que
décrit l'efficience faible au sens de Fama~\cite{fama1970}, la dépendance qui
subsiste étant trop ténue pour couvrir le coût de son exploitation. On retrouve ici, sur le marché
lui-même, la mise en garde formulée au premier test, un grand échantillon rend
significatif le moindre écart et il ne faut pas confondre significativité et
taille d'effet.

Le verdict du protocole s'éclaire alors. Si le passé du marché ne contient
presque rien d'exploitable, il est cohérent que des stratégies qui ne lisent
que ce passé n'aient rien trouvé. Le protocole n'a donc pas été
inutilement sévère, il a simplement constaté une absence qui était
attendue.

% =====================================================================
\section{Positionnement dans la littérature financière}
\label{sec:b1-positionnement}

Les outils de ce chapitre ont été construits à partir des seuls concepts de
cours, sondages, analyse de la variance et séries temporelles, or ils
coïncident avec les standards de l'économétrie financière.

L'agrégation par mois, une statistique par période puis l'inférence sur la
série de ces statistiques, est exactement la démarche de Fama et
MacBeth~\cite{famamacbeth1973}, standard du domaine. La variance de long terme
estimée par le modèle autorégressif poursuit quant à elle le même objectif que
les erreurs-types robustes à l'autocorrélation de Newey et
West~\cite{neweywest1987}, dont elle est une version paramétrique, validée ici
par le test de Box et Pierce.

Cette convergence est un argument en faveur du protocole. Les méthodes n'ont
pas été choisies parce qu'elles étaient reconnues en finance, mais parce que
le raisonnement statistique y menait ; le fait de retrouver au bout du chemin
des outils que la discipline utilise déjà indique que ce raisonnement était le
bon.

% =====================================================================
\section{Synthèse du chapitre}
\label{sec:b1-synthese}
Ce chapitre a procédé par une escalade de prudence. On commence par le test
de Student, le point de départ le plus simple pour une telle analyse. Malgré
ses hypothèses simplificatrices, il suffit déjà à écarter huit des onze
stratégies étudiées.

On a ensuite dû prendre en compte la dépendance, restée jusque-là hors du
champ du test. En reprenant les résultats du chapitre~\ref{chap:donnees}, qui
avait mis en évidence trois canaux de dépendance se ramenant à deux portes,
la dépendance entre les trades d'un même titre d'une part, celle liée au
moment d'autre part, il a fallu accepter de perdre beaucoup d'information en
agrégeant par titre puis par mois.

Cette agrégation aurait suffi s'il n'était pas resté, en outre, la dépendance
entre mois voisins. Il a donc fallu appliquer une dernière correction, qui a
elle-même exigé de vérifier la stationnarité des séries. Au terme de cette
escalade, plus aucune stratégie ne subsiste en dehors du témoin. Le verdict
positif du premier test ne tenait donc que par une condition qu'on n'avait pas
vérifiée.

Ce résultat n'a rien de surprenant une fois qu'on a interrogé le marché
lui-même. Son passé contient bien un peu d'information sur son avenir, mais
tellement peu que les frais l'effacent. Des stratégies qui ne lisent que ce
passé ne pouvaient donc pas trouver grand-chose.

Le chapitre suivant expose les limites de ce protocole, qui pourra certainement
être perfectionné avec le temps.